\cotefile{13}
\scanpage{46}
% loc: heading, lines 1–2
\textcircled{$\alpha$}. \uline{Foncteurs pleinement fidèles sur catégories des motifs}, \uline{liés à théorie de Hodge.}

% loc: opening, lines 3–5
$S$ un $\mathbb{Q}$-schéma régulier\inserted{ connexe}, muni d'un point \unsure{géom} $\unsure{\xi|\zeta}:\operatorname{Spec}(\mathbb{C})\longrightarrow S$
Le foncteur $\xi^*:\underset{\text{\inserted{\unsure{motifs lisses sur} $S$}}}{\operatorname{Mot}(S)}\longrightarrow\operatorname{Mot}(\mathbb{C})$ est fidèle. On \unsure{regardera} le foncteur composé
% The hook below Mot(S) attaches the gloss to the source.
% loc: diagram below the opening
\diagram{\begin{tikzcd}[ampersand replacement=\&,column sep=2.1em]
\operatorname{Mot}(S) \arrow[r,"\xi^*"] \arrow[rrr,bend right=14,"{(T_{\mathrm B})_\xi}"'] \& \operatorname{Mot}(\mathbb C) \arrow[r,"{T_{\deleted{\unsure{\mathrm B}}\unsure{\mathrm{B+H}|\mathrm{B,H}}}}"] \& \operatorname{Hodge}(\mathbb Q,\mathbb C) \arrow[r] \& \operatorname{Modf}(\mathbb Q)
\end{tikzcd}}

% loc: left of item 1
% slanted
\margin{réseau et vectoriel filtré}
% loc: item 1, first line
\textcircled{1} \inserted{(Via Hodge)} Supposons $\xi$ localisé au $\mathrm{p}^{\mathrm t}$ générique $\unsure{y|\eta}$ de $S$. \deleted{Le \unsure{fonct}} \textsuperscript{\textcircled{a}} \unsure{A|Á|À} tout motif $M$ sur $S$, on associe le triple
% loc: first triple
\[
\bigl\{T_{\mathrm{DR}}(M_y),T_{\mathrm B}(M_\xi),\varphi_M\bigr\}
\]
% loc: first comparison, line below triple
où $\varphi_M$ est l'isom. can. $\varphi_M:T_{\mathrm B}(M_\xi)\otimes_{\mathbb Q}\mathbb C\simeq T_{\mathrm{DR}}(M_y)\otimes_{k(y)}\mathbb C$.

% loc: slanted question by the two comparisons
% slanted
\margin{Le \unsure{foncteur} reste-t-il pl. fidèle si on laisse tomber la \unsure{filtration} de $T_{\mathrm{DR}}$ ?? id pour bigr. de $T_{\mathrm{Hdg}}$ !!}
% loc: filtration paragraph
On considère $T_{\mathrm{DR}}(M_y)$ comme \uline{filtré}, la filtration étant telle que $\varphi_M$ la transforme en une filtration \uline{admissible} du \unsure{complexifié} de $T_{\mathrm B}(M_\xi)$ [i.e. définissant, avec la filtration complexe conjuguée, une \unsure{bigraduation}].

% loc: second functor, lines above second triple
Le foncteur \inserted{en $M$} ainsi obtenu est pleinement fidèle (moyennant conj. de Hodge simple).\textsuperscript{\textcircled{b}} Même résultat si on considère
% loc: second triple
\[
\bigl\{T_{\mathrm{Hdg}}(M),T_{\mathrm B}(M_\xi),\Psi_M\bigr\}
\]
% loc: left of second triple
% slanted
\margin{réseau et vectoriel bigradué}
% loc: final comparison
$T_{\mathrm{Hdg}}(M)$ étant considéré comme \uline{bigradué}, et $\Psi_M$ étant l'isomorphisme canonique $\Psi_M:T_{\deleted{\unsure{\mathrm{DR}|\mathrm B}}\mathrm B}(M_\xi)\otimes_{\deleted{\unsure{k|\mathbb Q}}\mathbb Q}\mathbb C\simeq T_{\unsure{\mathrm{Hgg}|\mathrm{Hdg}}}(M)\otimes_{k(y)}\mathbb C$.
% loc: paper watermark, read sideways and compared across sheets
% printed form
% paper watermark, visible letters in the direction of the paper lettering
\par\noindent\textsc{trong}\hfill\textsc{Extra}\\
\unsure{ry|y} 1808\hfill\unsure{Maunu|Mauno|Maun}
% archive sheet number: 45

\scanpage{47}
% loc: item 2, above first line
\textcircled{2} \inserted{(Via Tate)} \inserted{Soit $\pi=\pi_1(S,\unsure{\xi|\zeta})$, soit $\ell$ un \unsure{nb} premier\unsure{,|.}} \unsure{A|À} tout $M$, associons
% loc: opening pair
\[
\bigl\{T_{\mathrm B\deleted{\unsure{-\mathrm H|+\mathrm H}}}(M_\xi),\rho_{\ell\unsure{,}M}^{\unsure{\text{ét}|\ell}}\bigr\}
\]
% loc: upper left, beside the cancelled block
% slanted
\margin{$\unsure{\mathbb Q_\ell|\ell}$-Module galoisien -- réseau}
% loc: two cancelled lines beneath pair
\deleted{où $T_{\mathrm{B}\unsure{.|\text{-}}\mathrm H}(M_\xi)$ est la structure de \deleted{\unsure{Betti-}}Hodge associée à $M_\xi$, \unsure{(}$T_{\mathrm{B}\unsure{+|\text{-}}\mathrm H}(M)=\{T_{\mathrm B}(M_{\deleted{\unsure{\xi}}\xi})+\text{bigraduation de }T_{\mathrm B}(M_\xi)\otimes_{\mathbb Q}\mathbb C\}$,)} \deleted{et} où
% loc: representation beneath cancellation
\[
\unsure{\rho_{\ell,M}|\rho_M}:\pi\longrightarrow\operatorname{Aut}_{\mathbb Q_\ell}\bigl(T_{\mathrm B}(M_\xi)\otimes_{\mathbb Q}\mathbb Q_\ell\bigr).
\]

% loc: first functor paragraph
Le foncteur ainsi obtenu est pl. fidèle si $S$ est \unsure{ess.} de type fini sur $\mathbb Q$ \deleted{$\unsure{\mathbb Q}$}\inserted{\deleted{\unsure{un corps|une conj.} $k$ ($\unsure{\hookrightarrow|\longrightarrow}\mathbb C$ par $\xi$)}}\unsure{,} en admettant comme de juste la conj. de \deleted{Hodge}\inserted{Tate}, et utilisant \unsure{rés.} de Griffiths

% loc: complex base and topological group
Si $S$ est \unsure{loc} de t.f. sur $\mathbb C$, on peut aussi prendre $\pi'=\pi_1(S(\mathbb C),\xi)$, et $\rho_M$ au lieu de $\rho_\ell$
\[
\rho_M:\pi'=\pi_1(S(\unsure{\mathbb C|\mathbb Q}),\xi)\longrightarrow\operatorname{Aut}_{\mathbb Q}\bigl(T_{\mathrm B}(M_\xi)\bigr),
\]
% loc: left of replacement by B-H
% slanted
\margin{ici c'est conj. Hodge qu'on utilise !}
à condition de remplacer $T_{\mathrm B}$ par $T_{\mathrm{B-H}}$ (structure de Hodge !).

% loc: affine case
Si $S$ est p.ex. affine, \deleted{\unsure{et}} mais pas néc. de type fini sur un corps, on a un résultat analogue, avec $\pi$ remplacé par un système projectif de groupes fondamentaux de \unsure{Schémas|Schém.} de t.f. sur $\mathbb Q$ \unsure{....}.

% loc: item 3 and quadruple
3) À tout $M$, associons
\[
\bigl(T_{\mathrm{DR}}(M),\sigma,T_{\mathrm B}(M_\xi),\varphi\bigr)
\]
% loc: description of item 3 and inserted qualifications
où $T_{\mathrm{DR}}(M)$ est la coh. de D.R. de $M$ sur $S$, considéré comme vectoriel \uline{filtré}\inserted{(il suffit de donner la filtration de $T_{\mathrm{DR}}(M)(\xi)$ !)}, $\sigma$ sa stratification de Gauss-Manin relativement à \deleted{$\unsure{\mathbb Q}$}\inserted{un corps de base fixé \uuline{Alg. clos} $k$ (p.ex. $k=\overline{\mathbb Q}$)}, $T_{\mathrm B}(M_\xi)$ la coh. de Betti,
% loc: lower slanted margin
% slanted
\margin{\unsure{on suppose} $S$ donné sur $k$\unsure{;} si $k=\mathbb Q$, il est inutile de mettre la filtration sur $T_{\mathrm{DR}}(M)$, cf plus bas.}
% loc: last comparison and its qualification
$\varphi$ l'isom. canonique $T_{\mathrm{DR}}(M)(\xi)\simeq T_{\mathrm B}(M_{\unsure{\xi|y}})\otimes_{\mathbb Q}\mathbb C$. Si $S$ est de t.f. sur $k$, on trouve ainsi un foncteur pl. fid. (moyennant conj. de Hodge\deleted{\unsure{§|?}}, et utilisant Griffiths).
% loc: paper watermark, read sideways and compared across sheets
% printed form
% paper watermark, visible letters in the direction of the paper lettering
\par\noindent\textsc{Extra Strong}\\
\unsure{Maunoury|Maunouy} 1808
% archive sheet number: 46

\scanpage{48}
% loc: item 4, heading and opening
\textcircled{4} \inserted{``\uline{Double réseau}''} Supposons $S=\operatorname{Spec}k$, $k\subset\mathbb C$ extension \uline{algébrique} de $\mathbb Q$. Considérons le foncteur qui à $M$ \unsure{associé|associe}
% loc: first labelled display
\begingroup\makeatletter\tagsleft@true\makeatother
\[
\bigl\{T_{\mathrm{DR}}(M),T_{\mathrm B}(M),\varphi_M\bigr\} \tag*{(4.1)}
\]
\endgroup
% loc: first comparison and question
où $\varphi:T_{\mathrm B}(M)\otimes_{\mathbb Q}\mathbb C\simeq T_{\mathrm{DR}}(M)\otimes_k\mathbb C$ ; ici, on ne considère pas $T_{\mathrm{DR}}(M)$ comme filtré, \deleted{\unsure{et}} ni $T_{\mathrm B}(M)$ muni d'une structure de Hodge. Ce foncteur est-il pl. fidèle ??

% loc: struck start and expression below it
\deleted{Il revient au $=$ de dire que}
\[
\deleted{T_{\unsure{\mathrm B|\mathrm R}}(M)\cap T_{\mathrm{DR}}(M)}
\]
% loc: reduction to k with overbar
Il \deleted{est} \unsure{aussitôt} ramené au cas où $k=\overline{\mathbb Q}$. La question est alors équivalente à celle si le foncteur associant à $M$ le système
% loc: second labelled display
\begingroup\makeatletter\tagsleft@true\makeatother
\[
\bigl\{T_{\mathrm B}(M),\sigma\bigr\} \tag*{(4.2)}
\]
\endgroup
% loc: connexion paragraph
où $\sigma$ est la connexion canonique sur $T_{\mathrm B}(M)\otimes_{\mathbb Q}\mathbb C$, est pleinement fidèle : en effet, $\sigma$ et $T_{\mathrm{DR}}(M)$\inserted{$\subset T_{\mathrm B}(M)\otimes_{\mathbb Q}\mathbb C$\unsure{,}} se déterminent mutuellement [$T_{\mathrm{DR}}(M)$ est l'\unsure{ens.} des éléments \inserted{de $T_{\mathrm B}(M)\otimes_{\unsure{k|\mathbb Q}}\mathbb C$} horizontaux pour $\sigma$ \ldots].
% loc: group description
Une autre façon de définir cette donnée, \inserted{$T_{\mathrm{DR}}\subset T_{\mathrm B}\otimes\mathbb C$,} quand on introduit le groupe de Galois $G=\underline{\operatorname{Aut}}_{\otimes}(T_{\mathrm B}\deleted{\unsure{M}})$ \deleted{\unsure{$\xi$} que \unsure{le}} (\unsure{$\mathcal M$|$M$} \unsure{sous-}$\otimes$-catégorie de type fini de la catégorie $\operatorname{Mot}(k)$) est la donnée de $\dot g\in\unsure{G|\mathbb G}(\mathbb C)/G(\overline{\mathbb Q})$ ; si $V$ est un $G$-module (correspondant à un motif $\deleted{\unsure{V}}M$), alors
% loc: short note in left margin
\margin{\uline{NB} $G$ connexe}
% loc: lattice and inclusion
$T_{\mathrm{DR}}(M)$ n'est autre que le $\overline{\mathbb Q}$-réseau $g(V_{\overline{\mathbb Q}})\subset V\otimes_{\mathbb Q}\mathbb C$. Donc la question est simplement si pour tout $V$, l'inclusion
\[
V^G\subset V\cap g(\overline V)
\]
% loc: last lines
est une identité. [On peut supposer que $V$ est un module simple. [S'il est de dim. 1, c'est une puissance de Tate, et ça marche].] \deleted{Il revient au $=$ de}
% loc: line of dashes under the final cancellation
\par\hfill - - - - - - - - - -
% loc: paper watermark, read sideways and compared across sheets
% printed form
% paper watermark, visible letters in the direction of the paper lettering
\par\noindent\textsc{Extra Strong}\\
\unsure{Maunoury|Maunouy} 1808
% archive sheet number: 47

\scanpage{49}
% loc: cancelled continuation at head
\deleted{regarder $V$ seulement (où $\overline V=V\otimes_{\mathbb Q}\overline{\mathbb Q}$)}
\[
\deleted{\overline V^{G}\subset\overline V\cap g(\overline V)}
\]
\deleted{et de prouver que c'est une identité [car $\overline V^G\cap V=V^G$].}

% loc: first equivalence
Or la relation $x\in\overset{\deleted{\unsure{-}}}{V}\cap g(\overline V)$ équivaut :
\begingroup\makeatletter\tagsleft@true\makeatother
\[
x\in\overset{\deleted{-}}{V},\quad g^{-1}x\in\overline V \tag*{\unsure{ds|$\alpha$)}}
\]
\endgroup
% loc: locus and second equivalence
et implique évidemment, si $\unsure{\overline T|\overline J}$\inserted{$\subset\overline G$} est \unsure{le|la} \deleted{\unsure{déclinaison|clôture} de} \inserted{lieu de} $g$ sur $G_{\overline{\mathbb Q}}$
\[
x\in\overline V,\quad\deleted{\unsure{g^{-1}x}}\quad t\deleted{\unsure{^{-1}}}s^{-1}\in G_x(\mathbb C)\quad\text{pour }t,s\in\overline T(\mathbb C),
\]
et $=$ y équivaut.
% loc: stabilizer paragraph
En d'autres termes, si $G_x$ est le stabilisateur de \deleted{$\unsure{G}$}$x$, qui est un ss-groupe algébrique de $G$, \deleted{\unsure{on}} \inserted{la relation $x\in V\cap g(\unsure{\overline W|\overline V})$ équi-}vaut :
\[
\deleted{\unsure{g}}\overline R\unsure{\in|\subset}\overline{G_x}\deleted{\unsure{(\mathbb C)}},
\]
où $\overline R$ est l'adhérence de l'im. $\overline T\times\overline T\longrightarrow\overline G$, $(s,t)\longmapsto s^{-1}t$.

% loc: image subgroup and equality question
Or soit $\widehat R$ l'image de $\overline R$ dans $G$, $H$ le sous-gr. qu'elle engendre, $H\subset G$. \deleted{Donc \unsure{la} propriété souhaitée équivaut à quelle \unsure{?}} La relation ds équivaut donc à $\deleted{\unsure{R}}H\subset G_x$ i.e. $H$ invarie $x$.
% loc: left insertion at opening of proposition
\inserted{La propriété \unsure{qu'on|qui} \unsure{veut est|nous est} celle-ci :}
% loc: heavily cancelled proposition and replacement
Pour tt $G$-module $V$, \deleted{et tt $x\in V$, tel que le stabilisateur $G_x$ contienne $\unsure{R}$, \inserted{i.e. tel que $R$ invarie $x$,} on a que $x\in V^G$ i.e. $G_x=G$.}
\deleted{Il suffit pour \unsure{ceci} que \unsure{le} \illegible{1} soit le ss-gpe engendré par $R$, soit égal à $G$. C'est aussi \unsure{nécessaire}} que les invariants de $H$ et de $G$ (\deleted{pour} \unsure{tt} repr. de $G$) soient les $=$\unsure{s}. Il suffit que l'on ait $H=G$, mais cette condition n'est pas nécessaire à priori, car un $k$-groupe de Borel de $G$ a cette propriété.
% loc: NB at left of cancelled lines
\margin{NB $\deleted{\unsure{\overline T}}\overline R,\overline H$ ne dépendent pas du choix de $g$ mod $G(\overline{\mathbb Q})$, \deleted{\unsure{à un près|a priori}}}
% loc: sloping note above 5
% slanted
\margin{\uline{Regarder} motifs \uline{associés} à \uline{courbes elliptiques}}
% loc: item 5, first half
5) Bien entendu, sans $=$ supposer que $k$ est alg. sur $\mathbb Q$, on peut regarder le foncteur (4.1). Mais il n'est plus en général pl. fidèle, $=$ si $k$ est de t.f. sur $\mathbb Q$, car pour $\unsure{\mathcal M|M}$ \inserted{$\otimes$-cat.} fixée de type fini $\subset\operatorname{Mot}(k)$, à condition d'agrandir $k$, on trouve que \inserted{les} $T_{\mathrm{DR}}(M)$ \uline{contiennent} $T_{\mathrm B}(M)$, donc les hom. des structures mixtes sont simplt les hom entre les $T_{\mathrm B}(M)$ : idiot !
% loc: closing lines of item 5
Pour $\mathcal M$ fixé, on peut espérer tout au mieux avoir quelque résultat de \unsure{disjonction} de $T_{\mathrm B}$ et $T_{\mathrm{DR}}$ si $k$ est $\pm$ \unsure{modeste|modérée|modulaire} (p.ex avoir des \unsure{estimations} de degré de transcendance minimum).
% loc: paper watermark, read sideways and compared across sheets
% printed form
% paper watermark, visible letters in the direction of the paper lettering
\par\noindent\textsc{Extra Strong}\\
\unsure{Maunoury|Maunouy} 1808
% archive sheet number: 48

\scanpage{50}
% loc: continuation of item 5
Par contre, sans hypothèse sur $k$ sauf d'être \uline{alg. clos}, il semble encore raisonnable d'espérer que (4.2) soit tjrs pleinement fidèle.
% loc: left oval
% slanted
\margin{réseau $+$ connexion}
On est ramené pour ceci au cas $k=\mathbb C$\unsure{..|.} C'est une conséquence de la conjecture ``du double réseau'' de 4) $+$ Tate ou Hodge, \unsure{préc.|précisé} \unsure{à|:} 6)

% loc: item 6, main line and interlinear insertion
6) Si $S$ de type fini \inserted{sur \deleted{$\unsure{\mathbb Q}$} un corps de nbs}, on a le foncteur de 3), mais sans filtration sur $T_{\mathrm{DR}}(M)$. Il résultera\deleted{\unsure{it|t}}it de la conj. du double réseau \deleted{(qui, elle, \unsure{génér} \inserted{(de \unsure{3|?})}} que ce foncteur est pl. fidèle
% loc: closing NB
[\uline{NB} si $S$ est le spectre d'un corps de nb, cela se réduit \unsure{:|à} l'énoncé du double réseau\unsure{....})
% loc: paper watermark, read sideways and compared across sheets
% printed form
% paper watermark, visible letters in the direction of the paper lettering
\par\noindent\textsc{\unsure{trong|rong}}\hfill\textsc{Extra}\\
\unsure{ry|y} 1808\hfill\unsure{Mauno|Maun}
% archive sheet number: 49

\scanpage{51}
% loc: heading above the rule
Remarque sur la caractérisation des classes de coh. D.R. \deleted{\unsure{qu}}\inserted{\unsure{essentiellement}} \uline{algébriques.} par bigraduation ou par réalité \unsure{....}
\par\noindent\rule{\linewidth}{0.3pt}\par
% loc: opening and Hom insertion
$K\subset\mathbb C$ sous-corps, $M$ motif sur $K$, d'où $T_{\mathrm{DR}}(M)$. Considérons $T_{\mathrm{DR}}(M\unsure{\leftrightharpoons|4|{}^{G}})$\inserted{$=\operatorname{Hom}(1,M)\otimes_{\mathbb Q}K$}, on voit que c'est un sous-module de $T_{\mathrm{DR}}(M)$, \deleted{qui est}\inserted{et que}
% loc: clauses a and b
\par a) contenu dans $T_{\mathrm{DR}}(M)_{\mathbb C}^{\unsure{00|\circ\circ}}$.
\par b) \inserted{$\operatorname{Hom}(1,M)\otimes_{\mathbb Q}(K\cap\mathbb R)$ est} contenu dans $\bigl(T_{\mathrm{DR}}(M)_{\mathbb C}\bigr)_{\mathbb R}$ \deleted{[et où $\operatorname{Hom}(1,M)\otimes_{\mathbb Q}(K\cap\mathbb R)$ est contenu dans $\bigl(T_{\mathrm{DR}}(M)_{\mathbb C}\bigr)_{\mathbb R}$]}

% loc: question
A-t-on une réciproque ?

% loc: item 1, opening and group insertions
$1^\circ$) Pas pour \deleted{\unsure{a) tout seul, même si}}\inserted{a) et b) simultanément} $K=\mathbb Q$, si le groupe de galois motivique $G$ \inserted{pour la cat. de motifs $\unsure{\mathcal M|M}$ où on travaille} n'est pas connexe. En effet, si $G$ est \deleted{\unsure{discret}} fini [Cas des motifs effectifs de poids 0] la condition a) est vide, car $T_{\mathrm{DR}}(M)_{\mathbb C}^{00}=T_{\mathrm{DR}}(M)_{\mathbb C}$, \inserted{donc réc. de a) vraie ssi $G=(e)$.} d'autre part \deleted{b) est} la réciproque à b) est vraie ssi l'extension galoisienne $L$ de groupe $G$ de $\mathbb Q$ définie par $\mathcal M$ est \uline{tot. réelle} [on teste avec le $G$-module $M=\mathbb Q[\unsure{C|G}]$, qui donne $T_{\mathrm{DR}}(M)=L$ \ldots] ; donc dans ce cas on a une réciproque à a) et b).

% loc: connected group and embeddings
Il faut donc se borner à un $G$ connexe, ce qui suggère de regarder les réciproques pour le cas $K$ alg. clos, \inserted{p.ex. $K=\overline{\mathbb Q}$} On peut de plus renforcer la condition, en regardant \uline{toutes} les immersions $K\hookrightarrow\mathbb C$.

% loc: item 2, introductory lines
$2^\circ$) La réciproque à a) est fausse, \deleted{\unsure{même si}} ($=$ pour $\overline{\mathbb Q}$). il suffit de prendre un groupe de Galois motivique $G$ qui est un \uline{tore} de dim $\geq 3$. On voit qu'on peut trouver \inserted{\unsure{sur|hom}}
% loc: pair of arrows, target, and marginal additions
$j_1,j_2:{\unsure{\mathbb G_{\mathrm m\,K}|\mathrm G_{\mathrm m\,K}}}^{\unsure{p|\prime}}\rightrightarrows G'_{\deleted{\unsure{K}}\unsure{K|\mathbb C}}$\inserted{$(=\operatorname{ad}(P)$, car $G$ \unsure{commut.})} qui défnt \inserted{(pour tt $K\longrightarrow\mathbb C$)} la bigraduation de
% loc: struck margin beside the two arrows
\margin{\deleted{\uline{NB} $j_1$ et $j_2$ \unsure{proviennent} d'\unsure{un choix} de \unsure{deux} \unsure{immersions} \unsure{dans} $\mathbb C$, et \illegible{1}}}
% loc: scalar extensions and torsor
$\deleted{\unsure{G}}V_{\mathbb C}\simeq V'\otimes_K\mathbb C$ \quad [$V'=P\mathbin{\overset{G_K}{\deleted{\unsure{\times}}}}V$, $P\unsure{\simeq|=}\underline{\operatorname{Isom}}_{\otimes}(T_{\mathrm B}\otimes K,T_{\mathrm{DR}})$]
% loc: degree-zero representation
Donc une représentation \inserted{fidèle} de degré 0 de $G$ \unsure{donne|donnée}, $\deleted{\unsure{V^{G'}}}$ si $T$ est l'image de $(j_1,j_2):\mathbb G_{\mathrm m\,K}\times\mathbb G_{\mathrm m\,K}\unsure{-|\longrightarrow}G'$, un $V'$ pour lequel
% loc: final equality and inequality
${V'}^{T}\inserted{(=T_{\mathrm{DR}}\cap(T_{\mathrm{DR}}\deleted{\unsure{)}}{}_{\mathbb C}^{00})}$
\inserted{\deleted{\unsure{quelque}}} est $\ne {V'}^{G'}$.
% loc: paper watermark, read sideways and compared across sheets
% printed form
% paper watermark, visible letters in the direction of the paper lettering
\par\noindent\textsc{Extra Strong}\\
\unsure{Maunoury|Maunouy} 1808
% archive sheet number: 50

\scanpage{52}
% loc: item 3, three lines at the head
$3^\circ$) Reste la question d'une réciproque à b) si $K$ \unsure{Alg.|alg.} clos, en prenant (\inserted{du moins} si $K\ne\overline{\mathbb Q}$) \uline{toutes} les imm. $K\longrightarrow\mathbb C$, et en \unsure{conjuguant} aussi au besoin avec la condition a).
% loc: paper watermark, read sideways and compared across sheets
% printed form
% paper watermark, visible letters in the direction of the paper lettering
\par\noindent\textsc{\unsure{ra|a} Strong}\hfill\textsc{Ex}\\
\unsure{aunoury|unoury|noury} 1808\hfill\unsure{M|A}
% archive sheet number: 51
